\honest{Feeling Rational}{Eliezer Yudkowsky}{2007}

Since curiosity is an emotion, I expect some people to object to counting it as part of rationality. \nb{``Why Truth?'', posted five months earlier, had named curiosity the most foundational reason to seek truth.}

Many people believe that rationality opposes all emotion. I can find no theorem of probability that says I should look ice-cold. I suspect such people are really contrasting ``System 1'' and ``System 2,'' fast perceptual judgments and slow deliberate ones. Each kind can serve the truth or defeat it.

Here is my test. I label an emotion ``not rational'' if it rests on mistaken beliefs, or rather on habits of thought that produce mistakes. I label it rational if correct beliefs or sound thinking evoke it. I quote a passage without saying where it comes from. If an iron comes toward your face and you think it hot when it is cool, the Way opposes your fear. If you think it cool when it is hot, the Way opposes your calm. \nb{It comes from Yudkowsky's own ``Twelve Virtues of Rationality.''} On this test, calm is an emotion like any other and not the rational default.

Our emotions ``arise from our models of reality.'' If I dream that my dead brother has been found alive, I am happy. When I wake, I am sad, and the sadness is rational. \nb{The post gives no evidence for the general claim. The dream is the case in which feeling follows belief at once. The post does not discuss feelings that stay after a belief has changed.}

If you believe a goblin tied your shoelaces together, that is a belief about the world: anything that can tie real laces must be real too. Anger at the goblin, I grant, is ``not just'' about the world. A calm person would expect the same knotted laces, since what happens in the closet does not depend on your mood. But the anger rests on the belief, so the standard of rationality spreads to it. \nb{The standard reaches only the belief under the anger. Whether the anger fits its cause, in kind or in size, the post does not judge.}

Becoming more rational can make feelings stronger, if you had been denying facts to escape them. I recall not knowing whether strong feeling was allowed. Since the days of Socrates at least, the way to seem cultured has been never to let anyone see you care strongly about anything. I cite nothing. \nb{Some periods prized the display of feeling. In the eighteenth-century novel of sensibility, the ability to display feelings was thought to show character.} People give me strange looks when they see how much I care. ``It's not the unusual subject, I think.'' They are not used to adults who visibly care about anything.

Now I know there is nothing wrong with feeling strongly. To Hodgell's ``That which can be destroyed by the truth should be'' I add my own rule: ``That which the truth nourishes should thrive.'' When something good happens, I am happy without wondering whether I may be. When something terrible happens, I do not flee my sadness by looking for ``fake consolations and false silver linings.'' \nb{The post links here to Yudkowsky's 2004 essay on the death of Yehuda Yudkowsky, the author's younger brother.} I picture the deaths and the hopes of humankind, and I admit that I am ``still embarrassed to confess strong emotions.'' ``But I know, now, that it is rational to feel.'' \nb{That follows from the post's own test: if a feeling is rational when its beliefs are true, then feeling is rational whenever they are.}
